% GENERATED FILE. Edit canonical lessons, then run npm run generate:books.
% canonical-source-hash: fnv1a64:7b079140449732d5
% canonical-lessons: KA-C539-edurali, KA-C539-nayaka, KA-C539-samabala, KA-C539-munnade, KA-C539-senasu, KA-R539-review-533-539

\chapter{Rivals}
\label{ch:ka-rivals}
\hlchaptermodality{539}

\begin{chapteropening}
\textbf{By the end of this chapter:} \emph{I can say an opponent, a captain, a draw, a lead and to contend.}

Complete the last lesson of chapter 539: I can say an opponent, a captain, a draw, a lead and to contend.
\end{chapteropening}

\section[\texorpdfstring{\kn{ಎದುರಾಳಿ} (edurāḷi)}{ಎದುರಾಳಿ (edurāḷi)}]{\kn{ಎದುರಾಳಿ} (edurāḷi) — an opponent}

\label{lesson:KA-C539-edurali}



\begin{quote}
Before the new one: say the Kannada for a tournament, then the Kannada for to defeat.
\end{quote}

\subsection*{You'll want to know: \kn{ಎದುರಾಳಿ}}
\textbf{\kn{ಎದುರಾಳಿ}} — \emph{edurāḷi} — \textquotedblleft{}an opponent\textquotedblright{}.

\begin{quote}
\kn{ಚದುರಂಗದಲ್ಲಿ} \kn{ನನ್ನ} \kn{ಎದುರಾಳಿ} \kn{ತುಂಬಾ} \kn{ಜಾಣ}.
\end{quote}

\begin{quote}
\emph{caduraṅgadalli nanna edurāḷi tumbā jāṇa.}
\end{quote}

\begin{quote}
My chess opponent is very clever.
\end{quote}

Plural: \kn{ಎದುರಾಳಿಗಳು}, \emph{edurāḷigaḷu}.

\kn{ಎದುರು}, facing + \kn{ಆಳು}, a man: the man facing you.

\begin{cousinweb}
\textbf{\kn{ಎದುರಾಳಿ}} (\emph{edurāḷi}) is a native Dravidian compound: \emph{edur}, \textquotedblleft{}that which is opposite\textquotedblright{} (DEDR 795), and \emph{āḷ}, \textquotedblleft{}a man, a soldier\textquotedblright{} (DEDR 399).
\end{cousinweb}

\begin{grammarlens}[title={What you've built}]
One more word for describing what you lived and felt.
\end{grammarlens}

\subsection*{Your turn}
\begin{itemize}
\raggedright
  \item \emph{Say it:} \emph{edurāḷi}
  \item \emph{Say it:} \emph{edurāḷi}, once more
  \item \emph{Say it:} \emph{caduraṅgadalli nanna edurāḷi tumbā jāṇa.} — My chess opponent is very clever.
  \item \emph{Recall:} say the Kannada for a tournament, then the Kannada for to defeat, then say \emph{edurāḷi} again
\end{itemize}

\begin{tcolorbox}[breakable,colback=teal!4,colframe=teal!35!black,title={Before you move on}]
What does \kn{ಎದುರಾಳಿ} mean? (\textquotedblleft{}An opponent\textquotedblright{}.) Say it once more.
\end{tcolorbox}
\section[\texorpdfstring{\kn{ನಾಯಕ} (nāyaka)}{ನಾಯಕ (nāyaka)}]{\kn{ನಾಯಕ} (nāyaka) — a captain; a leader}

\label{lesson:KA-C539-nayaka}



\begin{quote}
Before the new one: say the Kannada for to defeat, then the Kannada for an opponent.
\end{quote}

\subsection*{You'll want to know: \kn{ನಾಯಕ}}
\textbf{\kn{ನಾಯಕ}} — \emph{nāyaka} — \textquotedblleft{}a captain; a leader\textquotedblright{}.

\begin{quote}
\kn{ಅವನು} \kn{ನಮ್ಮ} \kn{ಕ್ರಿಕೆಟ್} \kn{ತಂಡದ} \kn{ನಾಯಕ}.
\end{quote}

\begin{quote}
\emph{avanu namma krikeṭ taṇḍada nāyaka.}
\end{quote}

\begin{quote}
He is the captain of our cricket team.
\end{quote}

Also the hero of a film; \kn{ನಾಯಕಿ}, a heroine.

\begin{cousinweb}
\textbf{\kn{ನಾಯಕ}} (\emph{nāyaka}) is borrowed from Sanskrit \emph{nāyaka}, \textquotedblleft{}a guide, leader, chief\textquotedblright{} (Monier-Williams), from \emph{nī}, \textquotedblleft{}to lead\textquotedblright{}.
\end{cousinweb}

\begin{grammarlens}[title={What you've built}]
One more word for describing what you lived and felt.
\end{grammarlens}

\subsection*{Your turn}
\begin{itemize}
\raggedright
  \item \emph{Say it:} \emph{nāyaka}
  \item \emph{Say it:} \emph{nāyaka}, once more
  \item \emph{Say it:} \emph{avanu namma krikeṭ taṇḍada nāyaka.} — He is the captain of our cricket team.
  \item \emph{Recall:} say the Kannada for to defeat, then the Kannada for an opponent, then say \emph{nāyaka} again
\end{itemize}

\begin{tcolorbox}[breakable,colback=teal!4,colframe=teal!35!black,title={Before you move on}]
What does \kn{ನಾಯಕ} mean? (\textquotedblleft{}A captain; a leader\textquotedblright{}.) Say it once more.
\end{tcolorbox}
\section[\texorpdfstring{\kn{ಸಮಬಲ} (samabala)}{ಸಮಬಲ (samabala)}]{\kn{ಸಮಬಲ} (samabala) — a draw; a tie}

\label{lesson:KA-C539-samabala}



\begin{quote}
Before the new one: say the Kannada for an opponent, then the Kannada for a captain.
\end{quote}

\subsection*{You'll want to know: \kn{ಸಮಬಲ}}
\textbf{\kn{ಸಮಬಲ}} — \emph{samabala} — \textquotedblleft{}a draw; a tie\textquotedblright{}.

\begin{quote}
\kn{ಎರಡು} \kn{ತಂಡಗಳ} \kn{ಪಂದ್ಯ} \kn{ಸಮಬಲದಲ್ಲಿ} \kn{ಮುಗಿಯಿತು}.
\end{quote}

\begin{quote}
\emph{eraḍu taṇḍagaḷa pandya samabaladalli mugiyitu.}
\end{quote}

\begin{quote}
The match between the two teams ended in a draw.
\end{quote}

\kn{ಸಮ}, equal + \kn{ಬಲ}, strength.

\begin{cousinweb}
\textbf{\kn{ಸಮಬಲ}} (\emph{samabala}) is a learned compound built from Sanskrit \emph{sama}, \textquotedblleft{}same, equal\textquotedblright{} (Monier-Williams), and \emph{bala}, \textquotedblleft{}power, strength\textquotedblright{} (CDIAL 9161).
\end{cousinweb}

\begin{grammarlens}[title={What you've built}]
One more word for describing what you lived and felt.
\end{grammarlens}

\subsection*{Your turn}
\begin{itemize}
\raggedright
  \item \emph{Say it:} \emph{samabala}
  \item \emph{Say it:} \emph{samabala}, once more
  \item \emph{Say it:} \emph{eraḍu taṇḍagaḷa pandya samabaladalli mugiyitu.} — The match between the two teams ended in a draw.
  \item \emph{Recall:} say the Kannada for an opponent, then the Kannada for a captain, then say \emph{samabala} again
\end{itemize}

\begin{tcolorbox}[breakable,colback=teal!4,colframe=teal!35!black,title={Before you move on}]
What does \kn{ಸಮಬಲ} mean? (\textquotedblleft{}A draw; a tie\textquotedblright{}.) Say it once more.
\end{tcolorbox}
\section[\texorpdfstring{\kn{ಮುನ್ನಡೆ} (munnaḍe)}{ಮುನ್ನಡೆ (munnaḍe)}]{\kn{ಮುನ್ನಡೆ} (munnaḍe) — a lead; an advance}

\label{lesson:KA-C539-munnade}



\begin{quote}
Before the new one: say the Kannada for a captain, then the Kannada for a draw.
\end{quote}

\subsection*{You'll want to know: \kn{ಮುನ್ನಡೆ}}
\textbf{\kn{ಮುನ್ನಡೆ}} — \emph{munnaḍe} — \textquotedblleft{}a lead; an advance\textquotedblright{}.

\begin{quote}
\kn{ಮೊದಲ} \kn{ಅರ್ಧದಲ್ಲಿ} \kn{ನಮ್ಮ} \kn{ತಂಡ} \kn{ಮುನ್ನಡೆ} \kn{ಪಡೆಯಿತು}.
\end{quote}

\begin{quote}
\emph{modala ardhadalli namma taṇḍa munnaḍe paḍeyitu.}
\end{quote}

\begin{quote}
In the first half our team took the lead.
\end{quote}

\kn{ಮುನ್}, ahead + \kn{ನಡೆ}, going.

\begin{cousinweb}
\textbf{\kn{ಮುನ್ನಡೆ}} (\emph{munnaḍe}) is a native Dravidian compound: \emph{mun}, \textquotedblleft{}that which is in front\textquotedblright{} (DEDR 5020), and \emph{naḍe}, \textquotedblleft{}to walk, go\textquotedblright{} (DEDR 3582).
\end{cousinweb}

\begin{grammarlens}[title={What you've built}]
One more word for describing what you lived and felt.
\end{grammarlens}

\subsection*{Your turn}
\begin{itemize}
\raggedright
  \item \emph{Say it:} \emph{munnaḍe}
  \item \emph{Say it:} \emph{munnaḍe}, once more
  \item \emph{Say it:} \emph{modala ardhadalli namma taṇḍa munnaḍe paḍeyitu.} — In the first half our team took the lead.
  \item \emph{Recall:} say the Kannada for a captain, then the Kannada for a draw, then say \emph{munnaḍe} again
\end{itemize}

\begin{tcolorbox}[breakable,colback=teal!4,colframe=teal!35!black,title={Before you move on}]
What does \kn{ಮುನ್ನಡೆ} mean? (\textquotedblleft{}A lead; an advance\textquotedblright{}.) Say it once more.
\end{tcolorbox}
\section[\texorpdfstring{\kn{ಸೆಣಸು} (seṇasu)}{ಸೆಣಸು (seṇasu)}]{\kn{ಸೆಣಸು} (seṇasu) — to contend; to fight hard}

\label{lesson:KA-C539-senasu}



\begin{quote}
Before the new one: say the Kannada for a draw, then the Kannada for a lead.
\end{quote}

\subsection*{You'll want to know: \kn{ಸೆಣಸು}}
\textbf{\kn{ಸೆಣಸು}} — \emph{seṇasu} — \textquotedblleft{}to contend; to fight hard\textquotedblright{}.

\begin{quote}
\kn{ಎರಡು} \kn{ತಂಡಗಳು} \kn{ಕೊನೆಯವರೆಗೂ} \kn{ಸೆಣಸಿದವು}.
\end{quote}

\begin{quote}
\emph{eraḍu taṇḍagaḷu koneyavaregū seṇasidavu.}
\end{quote}

\begin{quote}
The two teams fought hard to the very end.
\end{quote}

Past (I …): \kn{ಸೆಣಸಿದೆ}, \emph{seṇaside}.

With whom: -\kn{ಜೊತೆ}. \kn{ಸೆಣಸಾಟ} is a hard fight.

\begin{cousinweb}
\textbf{\kn{ಸೆಣಸು}} (\emph{seṇasu}) is a native Kannada verb, \textquotedblleft{}to struggle, to contend\textquotedblright{}; its origin is not settled.
\end{cousinweb}

\begin{grammarlens}[title={What you've built}]
One more word for describing what you lived and felt.
\end{grammarlens}

\subsection*{Your turn}
\begin{itemize}
\raggedright
  \item \emph{Say it:} \emph{seṇasu}
  \item \emph{Say it:} \emph{seṇasu}, once more
  \item \emph{Say it:} \emph{eraḍu taṇḍagaḷu koneyavaregū seṇasidavu.} — The two teams fought hard to the very end.
  \item \emph{Recall:} say the Kannada for a draw, then the Kannada for a lead, then say \emph{seṇasu} again
\end{itemize}

\begin{tcolorbox}[breakable,colback=teal!4,colframe=teal!35!black,title={Before you move on}]
What does \kn{ಸೆಣಸು} mean? (\textquotedblleft{}To contend; to fight hard\textquotedblright{}.) Say it once more.
\end{tcolorbox}
\section[(dialogue)]{Looking back: from Handwork to Rivals}

\label{lesson:KA-R539-review-533-539}



\begin{quote}
Say the Kannada for embroidery and handicraft.
\end{quote}

\subsection*{Your turn}
Say these aloud:

\begin{itemize}
\raggedright
  \item embroidery, handicraft, to knit, to string beads, to adorn
  \item a melody, a rhythm, a musical note, to hum a tune, to play an instrument
  \item a musical instrument, a veena, a flute, a violin, to beat a drum
  \item a stage play, theatre, a mask, to perform, to act
  \item a poet, a poem, literature, a work of art, to compose
  \item a medal, an athlete, a winner, a tournament, to defeat
  \item an opponent, a captain, a draw, a lead, to contend
\end{itemize}

\begin{tcolorbox}[breakable,colback=teal!4,colframe=teal!35!black,title={Before you move on}]
Embroidery? (\textbf{\kn{ಕಸೂತಿ}}, \emph{kasūti}.) Handicraft? (\textbf{\kn{ಕರಕುಶಲ}}, \emph{karakuśala}.) To knit? (\textbf{\kn{ಹೆಣೆ}}, \emph{heṇe}.) To string beads? (\textbf{\kn{ಪೋಣಿಸು}}, \emph{pōṇisu}.) To adorn? (\textbf{\kn{ಸಿಂಗರಿಸು}}, \emph{siṅgarisu}.) A melody? (\textbf{\kn{ರಾಗ}}, \emph{rāga}.) A rhythm? (\textbf{\kn{ತಾಳ}}, \emph{tāḷa}.) A musical note? (\textbf{\kn{ಸ್ವರ}}, \emph{svara}.) To hum a tune? (\textbf{\kn{ಗುನುಗು}}, \emph{gunugu}.) To play an instrument? (\textbf{\kn{ನುಡಿಸು}}, \emph{nuḍisu}.) A musical instrument? (\textbf{\kn{ವಾದ್ಯ}}, \emph{vādya}.) A veena? (\textbf{\kn{ವೀಣೆ}}, \emph{vīṇe}.) A flute? (\textbf{\kn{ಕೊಳಲು}}, \emph{koḷalu}.) A violin? (\textbf{\kn{ಪಿಟೀಲು}}, \emph{piṭīlu}.) To beat a drum? (\textbf{\kn{ಬಾರಿಸು}}, \emph{bārisu}.) A stage play? (\textbf{\kn{ನಾಟಕ}}, \emph{nāṭaka}.) Theatre? (\textbf{\kn{ರಂಗಭೂಮಿ}}, \emph{raṅgabhūmi}.) A mask? (\textbf{\kn{ಮುಖವಾಡ}}, \emph{mukhavāḍa}.) To perform? (\textbf{\kn{ಪ್ರದರ್ಶಿಸು}}, \emph{pradarśisu}.) To act? (\textbf{\kn{ಅಭಿನಯಿಸು}}, \emph{abhinayisu}.) A poet? (\textbf{\kn{ಕವಿ}}, \emph{kavi}.) A poem? (\textbf{\kn{ಕವಿತೆ}}, \emph{kavite}.) Literature? (\textbf{\kn{ಸಾಹಿತ್ಯ}}, \emph{sāhitya}.) A work of art? (\textbf{\kn{ಕಲಾಕೃತಿ}}, \emph{kalākṛti}.) To compose? (\textbf{\kn{ರಚಿಸು}}, \emph{racisu}.) A medal? (\textbf{\kn{ಪದಕ}}, \emph{padaka}.) An athlete? (\textbf{\kn{ಕ್ರೀಡಾಪಟು}}, \emph{krīḍāpaṭu}.) A winner? (\textbf{\kn{ವಿಜೇತ}}, \emph{vijēta}.) A tournament? (\textbf{\kn{ಪಂದ್ಯಾವಳಿ}}, \emph{pandyāvaḷi}.) To defeat? (\textbf{\kn{ಸೋಲಿಸು}}, \emph{sōlisu}.) An opponent? (\textbf{\kn{ಎದುರಾಳಿ}}, \emph{edurāḷi}.) A captain? (\textbf{\kn{ನಾಯಕ}}, \emph{nāyaka}.) A draw? (\textbf{\kn{ಸಮಬಲ}}, \emph{samabala}.) A lead? (\textbf{\kn{ಮುನ್ನಡೆ}}, \emph{munnaḍe}.) To contend? (\textbf{\kn{ಸೆಣಸು}}, \emph{seṇasu}.)
\end{tcolorbox}
